\section{Self-Attention 是内容决定的加权读取}
\label{sec:14-Deep-Learning/05-Attention-and-Sequence-Models/01}

读两句话：

\begin{center}
"猫坐在垫子上，因为它很舒服。"
\end{center}

\begin{center}
"猫坐在垫子上，因为它很柔软。"
\end{center}

读到"它"的时候，你在第一句里自动把它关联到"猫"——舒服的通常是动物。在第二句里自动关联到"垫子"——柔软的通常是织物。两个句子的词序和句法结构完全相同，"它"的指向取决于词的含义，不取决于位置。这个判断你做得如此自然，以至于不会停下来想它有多难：机器看到的只是一串浮点数向量，没有"动物"和"织物"的概念，没有常识，只有数字。怎么让它也学会"看内容决定读哪里"？

最粗暴的方案：给每个位置分配固定的读取窗口。位置 5 永远读位置 1 到 4 的平均。对刚才的句子勉强能工作——"它"在位置 5，"猫"和"垫子"都在窗口内。换一句："张三告诉李四他通过了考试"——"他"指谁？固定窗口覆盖了所有候选，但没有选择机制。窗口把信息拌在一起，无法按内容区分该听谁。

另一个极端：所有位置共享同一个全局表示。这锅粥更稠——"猫"和"垫子"的信息被搅成糊状，位置 5 拿到的只是"句子涉及猫和垫子"，但不知道当前这个词具体指代哪个。

两种方案栽在同一个地方：把"该读谁"写死在了位置或平均值里，而真答案藏在内容中。注意力的洞见：把读取目标变成一个可学习的、内容依赖的计算——每个位置发出查询，与所有位置的内容做匹配，按匹配度取信息。

\subsection{Query、Key、Value 把匹配和取值拆成三个空间}

设序列 \(X\in\R^{T\times D}\)。单头自注意力用三组独立的线性投影把每个位置映射到三个不同的角色：

\[
Q = XW_Q,\qquad K = XW_K,\qquad V = XW_V,
\]

其中 \(W_Q,W_K\in\R^{D\times d_k}\)，\(W_V\in\R^{D\times d_v}\)。同一个 \(x_t\) 被投影三次：Query 是"我想找什么"，Key 是"我有什么标签"，Value 是"我能提供什么内容"。同一位置在扮演不同角色时使用不同特征——作为查询者时暴露寻找方向，作为被查询者时暴露匹配标签，作为信息源时暴露传递内容。

分数矩阵

\[
S = \frac{QK^{\trans}}{\sqrt{d_k}} \in \R^{T\times T}
\]

的元素 \(S_{ts}\) 衡量位置 \(t\) 的查询与位置 \(s\) 的键有多匹配——两个向量的方向越一致，点积越大。沿每一行对可读位置 \(s\) 做 Softmax：

\[
A_{ts} = \frac{\exp(S_{ts})}{\sum_{r=1}^{T}\exp(S_{tr})},\qquad
Z = AV.
\]

每个输出位置得到一个加权和：

\[
z_t = \sum_{s=1}^{T} A_{ts} v_s.
\]

权重非负且和为 1——所有 Value 的凸组合。Query-Key 匹配决定"读谁、读多少"，Value 决定"读到什么"。三组投影的分离让模型用一套特征计算相关性，用另一套特征传递信息，而不是让同一组数字同时承担匹配和内容两种角色。

用一个只有三个 token 的序列把整个过程走一遍。"猫 坐 垫子"，\(D=4\)，投影后 \(d_k=d_v=2\)。

\[
Q = \begin{bmatrix} 0.5 & 1.2 \\ -0.3 & 0.8 \\ 0.1 & -0.5 \end{bmatrix},\quad
K = \begin{bmatrix} 0.7 & 0.9 \\ -0.1 & 0.6 \\ 0.2 & -0.3 \end{bmatrix},\quad
V = \begin{bmatrix} 1.0 & 0.2 \\ 0.5 & 1.1 \\ 0.3 & 0.8 \end{bmatrix}.
\]

分数矩阵（先忽略 \(\sqrt{d_k}\) 缩放）：

\[
QK^{\trans} = \begin{bmatrix}
0.5\cdot 0.7+1.2\cdot 0.9 & 0.5\cdot(-0.1)+1.2\cdot 0.6 & 0.5\cdot 0.2+1.2\cdot(-0.3) \\
-0.3\cdot 0.7+0.8\cdot 0.9 & -0.3\cdot(-0.1)+0.8\cdot 0.6 & -0.3\cdot 0.2+0.8\cdot(-0.3) \\
0.1\cdot 0.7+(-0.5)\cdot 0.9 & 0.1\cdot(-0.1)+(-0.5)\cdot 0.6 & 0.1\cdot 0.2+(-0.5)\cdot(-0.3)
\end{bmatrix}.
\]

第一行"猫"作为 Query：与三个 Key 的点积分別是 1.43、0.67、-0.26。Softmax 后权重约 \([0.62, 0.29, 0.09]\)——"猫"更多地从自身读取，少量参考"坐"，几乎不看"垫子"。第二行"坐"作为 Query：点积 0.51、0.51、-0.30，权重约 \([0.45, 0.45, 0.10]\)——均衡地关注"猫"和自身。每一次 Query 的内容变了，匹配结果就全变了。

\(z_1 = 0.62\cdot[1.0,0.2] + 0.29\cdot[0.5,1.1] + 0.09\cdot[0.3,0.8] = [0.80, 0.51]\)。Query 通过匹配权重决定"读多少"，Value 决定"读到什么"。\(W_Q,W_K,W_V\) 三组投影让模型从同一个输入 \(X\) 中提取出"用于匹配的特征"和"用于传递的内容"两套不同的东西。

\subsection{\(\sqrt{d_k}\) 在控制 Softmax 的饱和点}

如果 \(q_t\) 和 \(k_s\) 的各坐标近似独立、均值为零、方差为一，点积

\[
q_t^{\trans}k_s = \sum_{j=1}^{d_k} q_{tj} k_{sj}
\]

的方差约为 \(d_k\)。每一项是两个独立零均值单位方差随机变量的乘积——方差 \(\E[q^2]\E[k^2] - \E[q]^2\E[k]^2 = 1\cdot 1 - 0 = 1\)。\(d_k\) 项之和的方差累积为 \(d_k\)，标准差 \(\sqrt{d_k}\)。

问题在 Softmax 的指数函数上。\(d_k=64\)：点积典型量级约 \(\sqrt{64}=8\)，\(\exp(8)\approx 3000\)，\(\exp(-8)\approx 0.0003\)，区分度适中。\(d_k=512\)：典型量级跳到 \(\sqrt{512}\approx 23\)，\(\exp(23)\approx 10^{10}\)。如果某个点积偏离到了 30 或 40，\(\exp(30)\approx 10^{13}\) 对 \(\exp(0)=1\) 形成碾压——Softmax 趋近 one-hot，梯度 \(\partial A/\partial S\) 在饱和区趋于零，参数更新几乎只流经一个位置。

除以 \(\sqrt{d_k}\) 把典型量级拉回 \(O(1)\)，与维度解耦，避免过早进入梯度消失的饱和区。

这个解释是初始化层面的。它假设 \(q,k\) 各分量独立且方差约为一——有 LayerNorm 或适当初始化时近似成立，训练中分布会漂移。\(\sqrt{d_k}\) 是工程上有效的锚点，不是对所有训练阶段都精确的定理。

\subsection{Softmax 沿哪条轴归一化就是模型的语义声明}

注意力必须沿键位置 \(s\)（列方向）归一化，使每个查询 \(t\) 的权重和为 1：

\[
\sum_s A_{ts} = 1.
\]

这在说："每个输出位置把注意力分配到各个输入位置"。若误沿查询轴（行方向）归一化，语义变成"每个输入位置被多少输出位置关注"——计算图仍然合法，程序不报错，loss 可能还在降，但你建模的根本不是你想建模的东西。

数值稳定的实现逐行减去最大分数：

\[
\softmax(s)_j = \frac{\exp(s_j - m)}{\sum_r \exp(s_r - m)},\qquad m = \max_r s_r.
\]

平移不改变输出——分子分母同时乘 \(\exp(-m)\)——但避免了 \(\exp(80)\) 在 float32 下溢出。极端情况：某行全部被 mask（所有合法位置加 \(-\infty\)），减去 \(m=-\infty\) 后分子分母都是 0，结果为 0/0。实现必须显式定义——返回均匀分布、全零向量还是报错——不能假装公式自动覆盖。

\subsection{Mask 在限制信息流向，不只是排除干扰}

因果注意力要求位置 \(t\) 只能读取 \(s\le t\)：

\[
\widetilde S_{ts} = \begin{cases}
S_{ts}, & s\le t,\\
-\infty, & s>t.
\end{cases}
\]

Softmax 后被 mask 位置的权重严格为零。Padding mask 防止真实 token 读取补齐空白位。两种 mask 都在定义信息流向图，服务于不同的边界。

训练语言模型时有一个隐蔽的陷阱。训练数据是完整句子，目标是在位置 \(t\) 预测 \(t+1\)。如果因果 mask 偏移了一格——位置 \(t\) 能读到 \(t+1\)——训练 loss 异常低，因为模型已经看到了答案。部署时逐 token 生成，未来答案尚未产生，模型从"能偷看答案"突然切换为"什么也看不到"，表现断崖式下跌。这不是泛化差距——部署时跑的是另一个函数。

在病情时间序列中预测下一小时的风险，同一规则意味着未来化验结果不能进入当前位置。若因果 mask 与目标错开一个位置，模型在训练和随机切分验证中获得异常漂亮的 loss——答案已沿注意力边泄漏进输入；逐步部署时未来信息消失，性能才突然崩塌。

因果 mask 不是在做一个数值优化技巧，而是在声明：预测那一刻，真实可用的信息究竟是哪些。

有限精度下，代替 \(-\infty\) 的大负数必须与 dtype 匹配。float16 表示下限约 \(-65504\)；填 float32 的 \(-10^9\) 转 float16 后可能截断，Softmax 后仍留非零概率。直接填负无穷的浮点表示则可能在后续运算中引入 NaN。mask 必须在 Softmax 前加到分数矩阵上，不能在 Softmax 后硬清零——后者破坏 \(\sum_s A_{ts}=1\)。

\subsection{多头让不同的读取规则并行运行}

为每个头 \(h\) 用独立投影：

\[
Z^{(h)} = \softmax\!\left(\frac{Q^{(h)}K^{(h)\trans}}{\sqrt{d_k}}\right)V^{(h)}.
\]

各头输出拼接后再映射：

\[
Z = \operatorname{Concat}(Z^{(1)},\ldots,Z^{(H)})W_O.
\]

总宽度 \(D=512\)，\(H=8\) 个头，每个头 \(d_k=D/H=64\)，总计算量与单头 \(d_k=512\) 大致相同。但每个头在 64 维子空间里独立匹配和读取——不同头可以学到不同的匹配规则：一个关心邻居，另一个关心远距离语义，第三个专管句法。

多头不保证每个头学到人类能命名的关系。它只是允许多个低维匹配机制并行存在。头数增加而总宽度固定时，单头维度下降：\(D=512\)，\(H=2\) 时每个头 \(d_k=256\)，匹配空间丰富；\(H=16\) 时每个头仅 \(d_k=32\)，在 32 维空间里 Query-Key 点积能编码的匹配精细度相当有限。

有些头可能是冗余的——剪掉对下游指标几乎无影响；另一些可能学到退化行为，比如所有 Query 都只关注第一个 token（attention sink）。头的数量不是越多越好，需要结合任务的信息路由需求选择。

\subsection{没有位置信息时，注意力只看内容集合}

设 \(P\) 是置换矩阵。无位置编码时，纯注意力满足

\[
\operatorname{Attn}(PX) = P\operatorname{Attn}(X).
\]

\(Q,K,V\) 同时被 \(P\) 置换，\(QK^{\trans}\) 变为 \(P(QK^{\trans})P^{\trans}\)——行和列同时置换，每个分数仍连接原来那对位置，只是行列顺序一起变了。行 Softmax 保留这种置换关系，加权求和后输出每一行仍对应原来查询位置。模型知道内容及其匹配强度，但不知道谁的编号是 1。

位置编码、相对位置偏置、RoPE 和因果 mask，各自以不同机制打破这种对称。正弦编码给每个位置加固定向量；相对偏置在分数矩阵上加距离依赖项；RoPE 把位置旋入 Q 和 K，使点积自然携带相对距离。因果 mask 提供偏序（不能看未来），但区分不了"相隔 1 步"和"相隔 100 步"——精确的距离建模仍需显式位置结构。

\subsection{\(T^2\) 是全局读取的标价}

分数矩阵 \(S\in\R^{T\times T}\) 的存储和矩阵乘法都是 \(O(T^2 d_k)\)。\(T=2048\)、\(d_k=64\)（GPT-2 规模）：约 4M 元素，16MB（float32）。\(T=8192\)：约 67M 元素，268MB。\(T=32768\)：约 1B 元素，4GB——单层注意力已超过多数 GPU 显存。此时分数矩阵远超线性投影和 FFN 的总开销，Transformer 的瓶颈就卡在注意力层。

自回归生成有一个重要优化：过去位置的 Key 和 Value 在生成新 token 时不变。缓存它们：

\[
K_{\text{cache}} = [K_{\text{old}};\,k_{\text{new}}],\qquad
V_{\text{cache}} = [V_{\text{old}};\,v_{\text{new}}].
\]

每步只需算新 token 的 Q 并与缓存 K 做点积。第 \(t\) 步分数计算从 \(O(t^2 d_k)\) 降到 \(O(t d_k)\)，但缓存占用 \(O(t(d_k+d_v))\) 随序列线性增长。\(L\) 层、\(H\) 个头的 Transformer，生成第 \(T\) 个 token 时 KV cache 总量为 \(L\times H\times T\times (d_k+d_v)\times 2\) 字节（fp16）。\(L=32\)、\(H=32\)、\(d_k=d_v=128\)、\(T=4096\)，约 2GB——消费级 GPU 已接近上限。

稀疏注意力、低秩近似和分块注意力试图削减 \(T^2\)，但限制了哪些位置能直接交互。每个位置只读窗口 \(w=256\) 的邻域：\(T=10000\) 时计算量从 \(10^8\) 降到 \(\approx 2.56\times 10^6\)——约 40 倍加速。但相隔 500 步的位置永远无法单层直接交互，必须依赖中间层逐跳传递。值不值得，取决于任务真正需要的依赖距离，不取决于"更快"这个孤立的数字。

\subsection{注意力图适合检查，不适合宣称因果}

\(A_{ts}\) 是当前层、当前头中 Value 的线性组合系数。"模型关注了第 \(s\) 个词"仅限于：在这个线性组合中，\(v_s\) 的权重较大。

但这无法独自支撑因果声明。后续层可以重新混合信息。Value 投影 \(W_V\) 决定传递内容——注意力权重高但 Value 接近零向量，有效信息仍然很少。同一功能可由不同权重+Value 组合实现，权重差异大但输出几乎相同。两个不同随机种子训出的模型在同一个例子上可能给出完全不同的注意力分布，却产生几乎一样的预测。

注意力图的实际用途是诊断。检查 mask 是否正确——上三角区出现非零亮度说明因果 mask 挂掉了。检查 collapsed attention——某个位置在所有头上垄断了绝大部分权重，通常意味着信息瓶颈或模型偷懒。检查读取范围——生成第 100 个 token 时权重全压在最近 5 个 token 上，说明理论上的全局注意力退化成了局部窗口。

用热力图宣称"模型因为这个特征做出预测"则远远不够。注意力权重的数值受参数初始化随机性、优化路径偶然性和 Value 投影对信息重新编码的三重干扰。权重高不一定意味因果重要，权重低也不一定意味无关——信息可能已通过中间层 Value 传递编码到了其他位置。结构正确、数值稳定、预测有效和解释可信，是四种需要使用不同工具分别建立的独立结论。注意力图是诊断工具，不是因果证明。
